% Requires amsmath, booktabs, and multirow.
\subsection{Application to electronic health record alerts for acute kidney injury}
\label{sec:trial3_application}

\paragraph{Trial objective and analysis population.}
We applied the proposed method to the double-blind, multicenter, parallel randomized
clinical trial of electronic health record alerts for acute kidney injury (AKI) reported
by \citet{wilson2021electronic}. The trial evaluated whether an automated electronic
pop-up alert, accompanied by a link to an AKI order set, improved mortality, dialysis, and
AKI progression among adult inpatients with AKI across six hospitals. Eligible patients
were randomized at the patient level within the electronic health record using simple
randomization with concealed allocation. The analysis included all $N=6030$ randomized
participants: $n_1=3059$ assigned to the alert and $n_0=2971$ assigned to usual care. We
coded the alert as $A_i=1$ and usual care as $A_i=0$, and used the known randomization
probability $e=1/2$ rather than the realized allocation fraction.

\paragraph{Fixed-horizon hierarchical endpoint.}
The original primary outcome was the occurrence of AKI progression, dialysis, or death
within 14 days of randomization. We used the same restriction horizon, $S=14$ days, but
organized the three components hierarchically by clinical severity. Specifically, we
defined $D_{1i}(S)=I(\text{death by day 14})$,
$D_{2i}(S)=I(\text{dialysis by day 14})$, and
$D_{3i}(S)=I(\text{AKI progression by day 14})$, giving
$\mathcal D_i(S)=\{D_{1i}(S),D_{2i}(S),D_{3i}(S)\}$. All components were adverse binary
outcomes, so 0 was favorable. All three component indicators were observed for every
randomized participant (0 missing values for each component); consequently, no outcome
imputation or outcome-based subject exclusion was required. Participants were compared
lexicographically in the order death, dialysis, and AKI progression; the participant with
value 0 at the first unequal component won, and a pair was tied if all three components
were equal. This was a status-based comparison at the common horizon $S$: no
component-specific event times, discharge times, or lengths of stay were used.

\paragraph{Analysis.}
We compared the unadjusted estimator $\bar\eta_m$ with the proposed linear direct
influence-function projection estimator $\widehat\eta_m^{\mathrm{lin}}$ described in
Section~\ref{sec:algorithm}. The adjustment set comprised the prespecified, fully observed
baseline variables hospital, age over 90 years, sex, ICU status at randomization,
Elixhauser comorbidity score, diabetes, and congestive heart failure. Categorical variables
were reference coded, and the resulting feature vector was centered without an intercept.
The adjusted elementary win and loss probabilities were transformed exactly to NB, LWR,
and LWO, with influence-function-based inference for both methods.

\paragraph{Results.}
Among the $n_1n_0=9{,}088{,}289$ treatment--control pairs, the Alert participant won
$1{,}637{,}306$ comparisons, lost $1{,}669{,}083$, and tied $5{,}781{,}900$, giving
\[
  (\bar\tau_w,\bar\tau_l,\bar\tau_t)
  =(0.1802,0.1837,0.6362).
\]
Table~\ref{tab:trial3_primary} shows little evidence of a treatment benefit on any win
measure. Linear adjustment reduced the analytical standard error by approximately $5\%$
and the estimated variance by approximately $9\%$, but did not alter the substantive
conclusion. On the natural scales, the unadjusted and adjusted win ratios were 0.981
(95\% CI, 0.875--1.100) and 0.957 (95\% CI, 0.859--1.068), respectively; the corresponding
win odds were 0.993 (95\% CI, 0.953--1.035) and 0.984 (95\% CI, 0.946--1.024).

\begin{table*}[t]
\centering
\small
\caption{Unadjusted and linearly adjusted estimates in the AKI alert trial. LWR and LWO are
reported on the log scale. PVR is the percent variance reduction relative to the
unadjusted estimator based on analytical standard errors.}
\label{tab:trial3_primary}
\setlength{\tabcolsep}{5pt}
\begin{tabular}{llrrrrrr}
\toprule
Measure & Method & Estimate & ASE & 95\% CI lower & 95\% CI upper & $p$-value & PVR, \% \\
\midrule
\multirow{2}{*}{NB}
  & $\bar\eta_m$ & -0.0035 & 0.0106 & -0.0243 & 0.0173 & 0.742 & -- \\
  & $\widehat\eta_m^{\mathrm{lin}}$ & -0.0079 & 0.0101 & -0.0277 & 0.0119 & 0.434 & 9.2 \\
\addlinespace[2pt]
\multirow{2}{*}{LWR}
  & $\bar\eta_m$ & -0.0192 & 0.0583 & -0.1335 & 0.0950 & 0.742 & -- \\
  & $\widehat\eta_m^{\mathrm{lin}}$ & -0.0434 & 0.0555 & -0.1522 & 0.0653 & 0.434 & 9.4 \\
\addlinespace[2pt]
\multirow{2}{*}{LWO}
  & $\bar\eta_m$ & -0.0070 & 0.0212 & -0.0486 & 0.0346 & 0.742 & -- \\
  & $\widehat\eta_m^{\mathrm{lin}}$ & -0.0158 & 0.0202 & -0.0554 & 0.0238 & 0.434 & 9.2 \\
\bottomrule
\end{tabular}
\end{table*}
